\section{Irreducible-block auxiliaries}

\begin{theorem}[Scalar fixed points in an irreducible unital block]%
    \label{thm:irreducible_unital_scalar_fixed_point}
    \lean{MPSTensor.fixed_eq_scalar_of_isIrreducibleTensor_unital}
    \leanok
    \uses{def:transfer_map, def:irreducible_tensor}
    If $A$ is irreducible and unital, $\E_A(\Id)=\Id$, and $\E_A(X)=X$, then $X$
    is a scalar multiple of $\Id$. This is the final fixed-point assertion in the
    proof of \cite[Theorem~4]{PerezGarcia2007Matrix}.
\end{theorem}

\begin{proof}\leanok
    Tensor irreducibility gives irreducibility of the completely positive
    transfer map. By Wolf's Theorem~6.6 for irreducible unital Kraus maps, every
    fixed point is scalar.
\end{proof}

\begin{theorem}[Unital gauge from a full-rank Perron--Frobenius eigenvector]%
    \label{thm:unital_gauge_of_pd_fixed_point}
    \lean{MPSTensor.spectralUnitalGauge_isUnital_of_transferMap_eigenvector}
    \lean{MPSTensor.unitalGauge_isUnital_of_transferMap_fixedPoint}
    \lean{MPSTensor.sameMPV_unitalGauge}
    \leanok
    \uses{def:transfer_map}
    Let $\rho$ be positive definite with $\E_A(\rho)=r\rho$ for some $r>0$. For
    $B^i := r^{-1/2}\rho^{-1/2}A^i\rho^{1/2}$,
    \begin{equation}\label{eq:canon_unital_gauge_identity}
        \sum_i B^i(B^i)^\dagger = \Id.
    \end{equation}
    If $r=1$, then the unscaled gauge $B^i := \rho^{-1/2}A^i\rho^{1/2}$ is unital
    and generates the same MPV family as $A$. This is the spectral-radius
    normalization and full-rank fixed-point gauge of
    \cite[Theorem~4]{PerezGarcia2007Matrix}.
\end{theorem}

\begin{proof}\leanok
    The positive square root of $\rho$ is invertible. Substituting
    $B^i=r^{-1/2}\rho^{-1/2}A^i\rho^{1/2}$ into $\sum_i B^i(B^i)^\dagger$ and
    using $\E_A(\rho)=r\rho$ gives
    (\ref{eq:canon_unital_gauge_identity}). For $r=1$ the change from $A$ to $B$
    is a similarity, so cyclicity of the trace gives equality of all MPV
    coefficients.
\end{proof}

\begin{lemma}[Range factorization of the support projection]
    \label{lem:support_proj_range_factorization}
    \lean{MPSTensor.exists_supportProj_eq_mul}
    \lean{MPSTensor.supportProj_mul_left_eq_self}
    \lean{MPSTensor.one_sub_supportProj_mul_mul_supportProj_eq_zero}
    \leanok
    \uses{def:support_proj}
    The support projection $\pi$ of a positive semidefinite matrix $\rho$
    factors as $\pi=\rho W$ for a spectral pseudo-inverse $W$, witnessing
    that the range of $\pi$ is the range of $\rho$.  For a rectangular
    matrix $Y$, the support projection $\pi$ of $YY^\dagger$ fixes $Y$,
    $\pi Y=Y$; and whenever $AY=YG$ for some matrix $G$ --- that is, $A$
    maps the range of $Y$ into itself --- the one-sided invariance
    $(\Id-\pi)A\pi=0$ holds.
\end{lemma}

\begin{proof}
    \leanok
    \uses{def:support_proj}
    In the spectral decomposition
    $\rho=U\diag(\lambda)U^\dagger$ take
    $W=U\diag(\lambda^{+})U^\dagger$, where $\lambda^{+}$ inverts the
    positive eigenvalues and vanishes on the kernel; then
    $\rho W=U\diag(\mathbf 1_{\lambda>0})U^\dagger=\pi$.  From
    $\pi(YY^\dagger)=YY^\dagger$ the matrix
    $((\Id-\pi)Y)((\Id-\pi)Y)^\dagger$ vanishes, hence $\pi Y=Y$.
    Finally $(\Id-\pi)A\pi=(\Id-\pi)AY(Y^\dagger W)
    =(\Id-\pi)Y\,G(Y^\dagger W)=0$.
\end{proof}

\begin{theorem}[Irreducibility under rescaled conjugation]
    \label{thm:irreducible_tensor_smul_conj}
    \lean{MPSTensor.isIrreducibleTensor_smul_conj}
    \leanok
    \uses{def:irreducible_tensor, def:support_proj,
        lem:support_proj_range_factorization}
    Let $B$ be an irreducible tensor, let $X$ be invertible, and let
    $c\ne 0$.  Then the rescaled conjugated tensor with letters
    $c\,X^{-1}B^vX$ is irreducible.  In particular the rescaled unital
    gauge of an irreducible tensor is irreducible; this is the
    normalization of \cite[lines~224--225]{Cirac2016MPDO_arXiv} applied to
    an irreducible corner.
\end{theorem}

\begin{proof}
    \leanok
    \uses{lem:support_proj_range_factorization}
    An invariant orthogonal projection $P\notin\{0,\Id\}$ of the
    conjugated tensor makes the columns of $Y=XP$ span a subspace invariant
    under every letter of $B$: $B^vY=YG_v$ with
    $G_v=X^{-1}B^vXP$.  The support projection of $YY^\dagger$ is then a
    nontrivial invariant orthogonal projection of $B$ by
    Lemma~\ref{lem:support_proj_range_factorization}: it is nonzero
    because $Y\ne 0$, and it is not the identity because any nonzero
    vector of the kernel of $P$ transports through $(X^\dagger)^{-1}$ to
    the kernel of $YY^\dagger$.  This contradicts irreducibility of $B$.
\end{proof}

\begin{theorem}[Unital Schwarz data for an irreducible corner]
    \label{thm:spectral_unital_gauge_schwarz_setup}
    \lean{MPSTensor.spectralUnitalGauge_schwarz_setup}
    \leanok
    \uses{def:irreducible_tensor, def:transfer_map,
        thm:unital_gauge_of_pd_fixed_point, thm:irreducible_tensor_smul_conj}
    Let $B$ be an irreducible tensor with a positive definite transfer
    eigenvector, $\E_B(\rho)=r\rho$ with $\rho>0$ and $r>0$.  Then the
    rescaled unital gauge $K^v=r^{-1/2}\rho^{-1/2}B^v\rho^{1/2}$ is a
    unital Kraus family, is an irreducible tensor with an irreducible
    transfer map, and its adjoint map has a positive definite fixed point.
    These are the hypotheses of the cyclic decomposition of the peripheral
    spectrum, \cite[Theorem~6.6]{Wolf2012Quantum}.
\end{theorem}

\begin{proof}
    \leanok
    \uses{thm:unital_gauge_of_pd_fixed_point, thm:irreducible_tensor_smul_conj,
        thm:psd_fp_pd_irred}
    Unitality is Theorem~\ref{thm:unital_gauge_of_pd_fixed_point};
    irreducibility is Theorem~\ref{thm:irreducible_tensor_smul_conj}.  The
    adjoint of a unital family is trace preserving, so its transfer map is
    a channel and has a positive semidefinite fixed point by the Ces\`aro
    argument; irreducibility of the adjoint map upgrades the fixed point to
    a positive definite one.
\end{proof}

\begin{lemma}[Eigenvalue transport to the unital gauge]
    \label{lem:eigenvalue_transport_unital_gauge}
    \lean{MPSTensor.hasEigenvalue_transferMap_spectralUnitalGauge}
    \leanok
    \uses{def:transfer_map}
    In the setting of
    Theorem~\ref{thm:spectral_unital_gauge_schwarz_setup}, every transfer
    eigenvalue $\mu$ of $B$ becomes the transfer eigenvalue $\mu/r$ of the
    rescaled unital gauge: the gauge conjugates the transfer map by the
    congruence with $\rho^{1/2}$ and divides it by $r$.  In particular the
    peripheral eigenvalues $e^{2\pi iq/p}$ of
    \cite[lines~225--230]{Cirac2016MPDO_arXiv} arise from the transfer
    eigenvalues of modulus $r$.
\end{lemma}

\begin{proof}
    \leanok
    If $\E_B(Z)=\mu Z$ then
    $Z'=\rho^{-1/2}Z\rho^{-1/2}$ is nonzero and satisfies
    $\E_K(Z')=(\mu/r)Z'$ by direct computation.
\end{proof}

\begin{theorem}[Unitary diagonal positive fixed point in TP gauge]\label{thm:form_II}
    \lean{MPSTensor.exists_unitary_diag_posDef_fixedPoint_of_TP_of_isIrreducibleTensor}
    \leanok
    \uses{def:irreducible_tensor, def:tp_kraus}
    Let $A$ be an irreducible MPS tensor with $\sum_i (A^i)^\dagger A^i = \Id$ and
    $D > 0$. Then there exist a unitary $U$ and a diagonal positive definite
    $\Lambda$ such that, for $B^i := U^\dagger A^i U$,
    \[
        \sum_i (B^i)^\dagger B^i = \Id,
        \qquad
        \E_B(\Lambda) = \Lambda.
    \]
\end{theorem}

\begin{proof}
    \leanok
    \uses{thm:psd_fp_pd_irred}
    The TP hypothesis makes $\E_A$ a channel with a nonzero PSD fixed point.
    Irreducibility makes $\E_A$ an irreducible CP map, and
    Theorem~\ref{thm:psd_fp_pd_irred} upgrades the fixed point to positive
    definite. The spectral theorem diagonalizes it by a unitary conjugation,
    which preserves the TP normalization.
\end{proof}

\begin{remark}[Burnside input for cumulative spanning]\leanok
    \uses{thm:irreducible_tensor_to_action, thm:burnside_matrix,
        lem:algspan_cumulative_span}
    The passage from irreducible tensors to full algebra spanning first passes
    through the irreducible action on $\C^D$. Burnside's theorem then says that
    the resulting irreducible subalgebra of $\MN{D}$ is all of $\MN{D}$. Once
    the algebra span is full, finite-dimensionality gives a single cumulative
    span that already equals $\MN{D}$.
\end{remark}

\begin{lemma}[The algebra span reaches a finite cumulative span]\label{lem:algspan_cumulative_span}
    \lean{MPSTensor.exists_cumulativeSpan_eq_top_of_algSpan_eq_top}
    \leanok
    \uses{def:alg_span, def:cumulative_span}
    If the generated unital algebra satisfies $\algspn(A) = \MN{D}$,
    then there exists $N$ such that $T_N(A) = \MN{D}$.
\end{lemma}

\begin{proof}
    \leanok
    Every element of $\algspn(A)$ lies in some cumulative span $T_n(A)$:
    the generators lie in $T_1(A)$, scalar multiples of the identity lie in
    $T_0(A)$, and the family $\{T_n(A)\}_{n \ge 0}$ is closed under addition
    and under multiplication with lengths adding.
    Since $\MN{D}$ is finite-dimensional, the ascending chain
    $T_0(A) \subseteq T_1(A) \subseteq \cdots$ stabilizes, so one can choose a
    single $N$ with $T_N(A) = \MN{D}$.
\end{proof}

\begin{lemma}[Irreducible action gives an irreducible tensor]\label{lem:irreducible_action_to_tensor}
    \lean{MPSTensor.isIrreducibleTensor_of_isIrreducibleAction}
    \leanok
    \uses{def:irreducible_action, def:irreducible_tensor}
    If the letters of $A$ act irreducibly on $\C^D$, then $A$ has no nontrivial
    invariant orthogonal projection.
\end{lemma}

\begin{proof}
    \leanok
    If $P$ were a nontrivial invariant orthogonal projection, then its range
    would be invariant under every letter of $A$.
    Irreducibility of the action forces this range to be either $0$ or all of
    $\C^D$, so $P$ must be $0$ or $\Id$, a contradiction.
\end{proof}

\section{Normal canonical form}\label{sec:normal_canonical_form}

This section fixes the two canonical-form predicates used downstream: the
block-injective form of Definition~\ref{def:is_canonical_form} and the spectral
\emph{normal} canonical form of Definition~\ref{def:is_normal_canonical_form}. It
then records how the self-overlap clause and gauge-phase equivalence follow from
primitivity, and closes with the periodicity-removal step that produces primitive
transfer maps.

Recall from Definition~\ref{def:mpv_overlap}
and Lemma~\ref{thm:overlap_trace} that the bilinear overlap is
\[
    O_{AB}(N) = \sum_{\sigma \in \{0,\ldots,d{-}1\}^N}
    V^{(N)}(A)_\sigma \overline{V^{(N)}(B)_\sigma} = \Tr(F_{AB}^N).
\]
This notation enters the self-overlap clause of the canonical-form predicate
below.

\begin{definition}[Canonical form conditions]\label{def:is_canonical_form}
    \lean{MPSTensor.IsCanonicalForm}
    \leanok
    \uses{def:injective}
    A collection of scaling factors $(\mu_k)_{k=1}^r$ and
    block tensors $(A_k)_{k=1}^r$ satisfies the
    \emph{canonical form conditions} if:
    \begin{enumerate}
        \item each $A_k$ is injective,
        \item each $A_k$ satisfies the TP normalization
              $\sum_i (A_k^i)^\dagger A_k^i = \Id$,
        \item the moduli are non-increasing:
              $|\mu_1| \ge |\mu_2| \ge \cdots \ge |\mu_r|$,
        \item all $\mu_k \neq 0$,
        \item every block bond dimension is positive: $D_k \ge 1$ for every~$k$,
        \item the self-overlap converges: $O_{A_k A_k}(N) \to 1$
              for each~$k$.
    \end{enumerate}
    Only this one-sided normalization is assumed;
    no unital condition is assumed here.
    Condition~(6) is a primitivity hypothesis:
    for a primitive channel
    (\cite[Theorem~6.7(3)]{Wolf2012Quantum}),
    $T^n(\rho) \to \rho_\infty$ for every initial state,
    which forces the self-overlap to converge to~$1$.
    See also~\cite[Theorem~6.8]{Wolf2012Quantum}
    for the completely-positive characterisation
    and Chapter~\ref{ch:spectral} for the transfer-operator gap
    proof.
\end{definition}

% rem:pgvwc07_vs_local_cf (deleted): Definition~\ref{def:is_canonical_form} is a
% local/conditional predicate assuming injectivity and the self-overlap limit;
% it is not the source canonical-form existence or uniqueness theorem of
% \cite[Theorem~4]{PerezGarcia2007Matrix}.

\begin{remark}[Normalization convention]\label{rem:tp_vs_unital}
    \cite[Theorem~4]{PerezGarcia2007Matrix} uses the unital gauge
    $\sum_i A_k^i (A_k^i)^\dagger = \Id$; here the blocks are in the dual TP gauge
    $\sum_i (A_k^i)^\dagger A_k^i = \Id$. A positive-definite fixed point of the
    adjoint transfer map provides the change of gauge between the two
    (Theorem~\ref{thm:left_canonical_gauge}).
\end{remark}

\begin{definition}[Normal canonical form conditions]\label{def:is_normal_canonical_form}
    \lean{MPSTensor.IsNormalCanonicalForm}\leanok
    \uses{def:irreducible_tensor, def:primitive_channel, def:tp_kraus}
    Scaling factors $(\mu_k)_{k=1}^r$ and block tensors $(A_k)_{k=1}^r$ satisfy
    the \emph{normal canonical form conditions} if:
    \begin{enumerate}
        \item each $A_k$ is irreducible,
        \item each $A_k$ satisfies $\sum_i (A_k^i)^\dagger A_k^i = \Id$,
        \item each $\E_{A_k}$ is primitive
              (equivalently, $\sigma_\partial(\E_{A_k}) = \{1\}$),
        \item $|\mu_1| \ge |\mu_2| \ge \cdots \ge |\mu_r|$,
        \item all $\mu_k \neq 0$,
        \item all bond dimensions are positive.
    \end{enumerate}
    Here ``normal'' is the spectral sense of~\cite{Cirac2016MPDO_arXiv}, which
    by~\cite[Proposition~3]{Sanz2010Quantum} is equivalent to the eventual
    block-injectivity formulation of Definition~\ref{def:normal}.
\end{definition}

% The weight moduli are non-increasing, not strictly decreasing; equal-modulus
% and repeated-copy sectors are compared by the sector-coefficient construction
% of Chapter~\ref{ch:bnt} (Definition~\ref{def:mpv_phase_class_data}, supplier
% Theorem~\ref{thm:paperbnt_supplier_prepared}), where a repeated sector
% contributes the power sum $c_N^{(j)} = \sum_q \mu_{j,q}^N$.

\begin{theorem}[Self-overlap is derived in normal canonical form]%
    \label{thm:ncf_overlap_tendsto_one}
    \lean{MPSTensor.IsNormalCanonicalForm.overlap_tendsto_one}\leanok
    \uses{def:is_normal_canonical_form, def:mpv_overlap}
    If $(\mu_k, A_k)$ satisfies the normal canonical form conditions, then
    $O_{A_k A_k}(N) \to 1$ for every block~$k$.
\end{theorem}

\begin{proof}\leanok
    \uses{def:is_normal_canonical_form, thm:primitive_compl_lt_one, thm:primitive_overlap}
    Fix a block~$k$. Irreducibility and the TP normalization give a nonzero
    positive semidefinite fixed point $\rho$ of $\E_{A_k}$ with fixed-point
    projector $P$; primitivity gives the complementary gap
    $\rho_{\spec}(\E_{A_k}-P) < 1$ (Theorem~\ref{thm:primitive_compl_lt_one}).
    Theorem~\ref{thm:primitive_overlap} then yields $O_{A_k A_k}(N) \to 1$.
\end{proof}

\begin{remark}
    See Theorem~\ref{thm:modulus_one_gauge_irr_tp} for the modulus-one
    eigenvalue rigidity statement for irreducible trace-preserving blocks.
\end{remark}

\begin{lemma}[Mixed-transfer spectral radius from unit-modulus overlap]%
    \label{lem:overlap_norm_one_implies_spectral_radius_ge_one}
    \lean{MPSTensor.mixedTransferSpectralRadius_ge_one_of_mpvOverlap_norm_tendsto_one}
    \leanok
    \uses{def:mpv_overlap}
    Let $A$ and $B$ have the same bond dimension. If
    $\|\braket{V^{(N)}(A)}{V^{(N)}(B)}\| \to 1$, then the mixed transfer map has
    spectral radius at least $1$.
\end{lemma}
\begin{proof}\leanok
    If the mixed transfer spectral radius were $<1$, the iterated mixed transfer
    map would tend to zero in operator norm, hence so would its trace; via
    $\Tr(F_{AB}^N) = \braket{V^{(N)}(A)}{V^{(N)}(B)}$ the overlap would tend
    to~$0$, contradicting the hypothesis.
\end{proof}

\begin{theorem}[Gauge-phase equivalence from unit-modulus overlap]%
    \label{thm:nt_overlap_norm_one_gauge}
    \lean{MPSTensor.gaugePhaseEquiv_of_overlap_norm_tendsto_one_of_irreducible_TP}
    \leanok
    \uses{def:irreducible_tensor, def:tp_kraus, def:gauge_phase_equiv,
        def:mpv_overlap}
    Let $A$ and $B$ be irreducible MPS tensors of the same bond dimension with
    $\sum_i (A^i)^\dagger A^i = \Id$ and $\sum_i (B^i)^\dagger B^i = \Id$. If
    $\|\braket{V^{(N)}(A)}{V^{(N)}(B)}\| \to 1$, then $A$ and~$B$ are gauge-phase
    equivalent.
    % Same-bond-dimension case of \cite[Lemma~equalMPS]{Cirac2016MPDO_arXiv};
    % proportionality of MPVs is not assumed.
\end{theorem}

\begin{proof}\leanok
    \uses{lem:overlap_norm_one_implies_spectral_radius_ge_one,
        thm:modulus_one_gauge_irr_tp}
    By Lemma~\ref{lem:overlap_norm_one_implies_spectral_radius_ge_one} the mixed
    transfer spectral radius is at least~$1$, and
    Theorem~\ref{thm:modulus_one_gauge_irr_tp} then gives gauge-phase equivalence.
\end{proof}

\subsection{Canonical form from primitivity}\label{sec:canonical_from_primitivity}

A normal block needs a primitive transfer map; an irreducible TP block may still
have peripheral spectrum larger than $\{1\}$, namely the roots of unity coming
from its period. Blocking by the least common multiple of those periods raises the
transfer map to a power whose peripheral spectrum is $\{1\}$, removing the
periodicity. In~\cite{Cirac2021Matrix} this is the passage from irreducibility to
primitivity by blocking to period one.
% rem:primitivity_reduces_assumptions / "two ways" (deleted): once
% peripheral-spectrum primitivity is known, the overlap-convergence hypothesis is
% redundant, derived from the transfer-map gap of Chapter~\ref{ch:spectral}; these
% are statements about blocks already in normal form, not the arbitrary-input
% existence theorem of \cite[Theorem~4]{PerezGarcia2007Matrix}.

\begin{theorem}[Blocking yields a peripheral-spectrum primitive transfer map]\label{thm:blocking_primitivity}
    \lean{MPSTensor.exists_blockTensor_isPrimitive_of_TP_of_isIrreducibleTensor}
    \leanok
    \uses{def:blocked_tensor, def:transfer_map, def:primitive_channel,
        def:tp_kraus, def:irreducible_tensor}
    Let $A$ be an irreducible MPS tensor with $D > 0$ and
    $\sum_i (A^i)^\dagger A^i = \Id$. Then there exists a blocking length $p > 0$
    such that $\sigma_\partial(\E_{A^{[p]}}) = \{1\}$.
\end{theorem}

\begin{proof}\leanok
    \uses{thm:peripheral_roots_irred_fp, thm:peripheral_finite,
        lem:common_power_eq_one, lem:peripheral_pow_singleton}
    Pass to the conjugate-transposed Kraus family $K_i := (A^i)^\dagger$, which is
    unital with irreducible Kraus map. The TP fixed-point theorem for $\E_A$
    provides a positive definite fixed point for the adjoint of $K$, so
    Theorem~\ref{thm:peripheral_roots_irred_fp} shows every peripheral eigenvalue
    is a root of unity. Since $\sigma_\partial(\E_A)$ is finite
    (Theorem~\ref{thm:peripheral_finite}),
    Lemma~\ref{lem:common_power_eq_one} gives a common exponent $p$ with
    $\mu^p = 1$ for all $\mu \in \sigma_\partial(\E_A)$, and
    Lemma~\ref{lem:peripheral_pow_singleton} gives
    $\sigma_\partial(\E_A^p) = \{1\}$. Blocking by $p$ takes the $p$th power of
    the transfer map, so $\E_{A^{[p]}}$ is primitive.
\end{proof}

\begin{remark}\label{rem:blocking_primitivity_appendixA}
    Theorem~\ref{thm:blocking_primitivity} is the periodicity-removal-by-blocking
    step of~\cite[Appendix~A]{Cirac2016MPDO_arXiv} for an irreducible block
    already in TP gauge.
\end{remark}

\section{Left-canonical and dual diagonalization of irreducible blocks}%
\label{sec:canonical_construction_1606}

Each irreducible block in trace-preserving gauge admits the PGVWC07 unital
orientation together with a diagonal positive-definite dual fixed point
(Theorem~\ref{thm:CFII_data}, the Perron--Frobenius gauge of
Chapter~\ref{ch:qpf}, and
Theorem~\ref{thm:irreducible_unital_scalar_fixed_point}); the full PGVWC07
canonical form is Theorem~\ref{thm:pgvwc07_ti_canonical_form}.

\begin{theorem}[Left-canonical normalization for irreducible TP blocks]%
    \label{thm:CFII_data}
    \lean{MPSTensor.exists_CFII_data_of_TP_of_isIrreducibleTensor}
    \leanok
    \uses{def:irreducible_tensor}
    Let $A$ be an irreducible MPS tensor in TP gauge
    ($\sum_i (A^i)^\dagger A^i = \Id$) with $D > 0$. Then there exist a unitary
    $U$ and a diagonal positive definite $\Lambda$ such that:
    \begin{enumerate}
        \item $B^i := U^\dagger A^i U$ is TP,
        \item $\E_B(\Lambda) = \Lambda$,
        \item $\Lambda$ is diagonal and positive definite,
        \item $A$ and $B$ generate the same MPV family.
    \end{enumerate}
    This is the left-canonical normalization step (Canonical Form~II)
    of~\cite[Appendix~A]{Cirac2016MPDO_arXiv}. The unitary conjugation is a
    second step inside TP gauge, performed only after TP normalization.
\end{theorem}

\begin{proof}\leanok
    \uses{thm:form_II, thm:samempv_unitary}
    Theorem~\ref{thm:form_II} provides the unitary and the diagonal PD fixed
    point, and unitary conjugation preserves trace-preservation and the MPV
    family (Theorem~\ref{thm:samempv_unitary}).
\end{proof}

\begin{theorem}[Dual diagonalization for irreducible unital blocks]%
    \label{thm:pgvwc07_dual_diag_unital}
    \lean{MPSTensor.exists_unitary_diag_posDef_adjointFixedPoint_of_unital_of_isIrreducibleTensor}
    \leanok
    \uses{def:irreducible_tensor}
    Let $A$ be an irreducible MPS tensor in unital gauge
    $\sum_i A^i(A^i)^\dagger=\Id$ with $D>0$. Then there exist a unitary $U$ and a
    diagonal positive definite $\Lambda$ such that, for $B^i:=U^\dagger A^iU$,
    \begin{equation}\label{eq:canon_dual_diag_unital_pair}
        \sum_i B^i(B^i)^\dagger=\Id,
        \qquad
        \sum_i (B^i)^\dagger\Lambda B^i=\Lambda,
    \end{equation}
    and $A$, $B$ generate the same MPV family. This is the dual-map fixed-point
    and unitary-diagonalization step of
    \cite[Theorem~4]{PerezGarcia2007Matrix}.
\end{theorem}

\begin{proof}\leanok
    \uses{thm:CFII_data, thm:samempv_unitary}
    Apply Theorem~\ref{thm:CFII_data} to the conjugate-transposed family
    $i\mapsto (A^i)^\dagger$: unitality of $A$ is trace preservation for this
    family, whose transfer map is the dual of $\E_A$. Translating back gives the
    equations~(\ref{eq:canon_dual_diag_unital_pair}).
\end{proof}

\begin{lemma}[Common scalar rescaling of block weights]%
    \label{lem:mpv_to_tensor_from_blocks_weight_mul_left}
    \lean{MPSTensor.mpv_toTensorFromBlocks_weight_mul_left}
    \leanok
    Multiplying every block weight in a block-diagonal MPS tensor by a scalar $c$
    multiplies every length-$N$ MPV coefficient by $c^N$.
\end{lemma}

\begin{proof}\leanok
    \uses{thm:mpv_decomposition}
    Expand the block-diagonal MPV coefficient as a finite sum over blocks and
    factor $c^N$ out of the sum.
\end{proof}

\begin{theorem}[Normal canonical form from a primitive block decomposition]%
    \label{thm:exists_normal_canonical_form}
    \lean{MPSTensor.exists_normalCanonicalForm_of_primitive_blockDecomp}
    \leanok
    \uses{def:is_normal_canonical_form, def:blocked_tensor}
    Suppose $A$ generates the same MPV family as $\bigoplus_k \mu_k A_k$, and that
    each $A_k$ is irreducible, satisfies $\sum_i (A_k^i)^\dagger A_k^i = \Id$, and
    has primitive transfer map, all $\mu_k \neq 0$, and all bond dimensions are
    positive. Then there exist a blocking length $p > 0$, weights $(\nu_k)$, and
    blocked tensors $(B_k)$ such that $A^{[p]}$ generates the same MPV family as
    $\bigoplus_k \nu_k B_k$, and $(\nu_k, B_k)$ satisfies the normal canonical
    form conditions. The weight moduli need not be distinct; equal-modulus blocks
    are allowed.
\end{theorem}

\begin{proof}\leanok
    \uses{def:is_normal_canonical_form}
    The given blocks are already primitive, so take $p = 1$. Reorder the blocks by
    non-increasing weight modulus, with ties allowed, and set $\nu_k := \mu_k$ and
    $B_k := A_k$ in that order. Irreducibility, TP normalization, primitivity,
    nonzero weights, and positive bond dimensions are preserved, so the reordered
    family satisfies Definition~\ref{def:is_normal_canonical_form} and generates
    the same MPV family as~$A$.
\end{proof}

\begin{remark}[Scope of Theorem~\ref{thm:exists_normal_canonical_form}]
    Theorem~\ref{thm:exists_normal_canonical_form} assumes all six block
    hypotheses in advance and produces the normal canonical form by reordering.
    It is not the arbitrary-input existence theorem
    of~\cite[Section~2.3]{Cirac2016MPDO_arXiv}; that primitive-block reduction,
    including its structural bond-dimension bound, is
    Theorem~\ref{thm:tp_primitive_blockdecomp} in
    Section~\ref{sec:reduction_to_normal_form}.
\end{remark}


\section{Support compression and period-removal auxiliaries}

\begin{theorem}[Supported corner compression with a support isometry]%
    \label{thm:supported_corner_compression_isometry}
    \lean{MPSTensor.exists_compressedTensor_of_supported_projection_with_letter_and_isometry}
    \leanok
    \uses{def:support_proj}
    Let $A$ be a tensor normalized on the support of an orthogonal projection
    $P$, so that $P A^i P=A^i$ and
    $\sum_i (A^i)^\dagger A^i=P$. Then there are a corner tensor $C$, a
    linear isomorphism
    \[
        \varphi:\MN{\dim C}\xrightarrow{\sim} P\MN{D}P,
    \]
    and a rectangular isometry $V:\C^{\dim C}\to\C^D$ such that
    \[
        V^\dagger V=\Id,\qquad VV^\dagger=P,\qquad
        \varphi(X)=VXV^\dagger.
    \]
    Moreover $\varphi(C^i)=A^i$ for every physical letter, and for every word
    $\sigma=(\sigma_1,\ldots,\sigma_N)$,
    \[
        V^{(N)}(C)_\sigma
        = \tr\!\bigl(P A^{\sigma_1}\cdots A^{\sigma_N}\bigr).
    \]
    \begin{proof}\leanok
        Diagonalize the projection $P$ by a unitary change of basis.  The
        inclusion of the support coordinates gives a rectangular isometry $V$
        with $V^\dagger V=\Id$ and $VV^\dagger=P$.  Compress the letters of
        $A$ to this support and let $\varphi$ be the corresponding corner
        isomorphism.  The support relation $P A^iP=A^i$ gives the
        corner-letter identity. For each word
        $\sigma=(\sigma_1,\ldots,\sigma_N)$ the MPV component is obtained from
        the trace computation
        \[
            V^{(N)}(C)_\sigma
            =\tr(C^{\sigma_1}\cdots C^{\sigma_N})
            =\tr(V^\dagger A^{\sigma_1}\cdots A^{\sigma_N}V)
            =\tr(PA^{\sigma_1}\cdots A^{\sigma_N}).
        \]
        The final equality uses cyclicity of the trace and $VV^\dagger=P$.
    \end{proof}
\end{theorem}

\begin{theorem}[Commuting projections give sector isometries]%
    \label{thm:commuting_projection_block_decomp_isometry}
    \lean{MPSTensor.exists_blockDecomp_of_commuting_projections_with_letter_and_isometry}
    \leanok
    \uses{thm:supported_corner_compression_isometry}
    Let a left-canonical tensor $A$ commute with a family of orthogonal
    projections $(P_k)$ summing to $\Id$. Then $A$ decomposes into compressed
    sector tensors $(C_k)$, with a unit-weight direct-sum MPV decomposition,
    corner isomorphisms $\varphi_k:\MN{\dim C_k}\to P_k\MN{D}P_k$, and
    rectangular isometries $V_k$ satisfying
    \[
        V_k^\dagger V_k=\Id,\qquad V_kV_k^\dagger=P_k,\qquad
        \varphi_k(X)=V_kXV_k^\dagger .
    \]
    The sector letters obey
    $\varphi_k(C_k^i)=P_k A^i P_k$.
    \begin{proof}\leanok
        Apply Theorem~\ref{thm:supported_corner_compression_isometry} to the
        compression of $A$ on each sector $P_k$.  Orthogonality and
        $\sum_k P_k=\Id$ give the direct-sum decomposition. For every word
        $\sigma=(\sigma_1,\ldots,\sigma_N)$ the partition of unity gives
        \[
            V^{(N)}(A)_\sigma
            =\tr\!\Bigl((\sum_k P_k)A^{\sigma_1}\cdots A^{\sigma_N}\Bigr)
            =\sum_k\tr(P_kA^{\sigma_1}\cdots A^{\sigma_N})
            =\sum_k V^{(N)}(C_k)_\sigma,
        \]
        which is the unit-weight matrix-product-vector identity.
    \end{proof}
\end{theorem}

\begin{theorem}[Adjoint-fixed projections give sector isometries]%
    \label{thm:adjoint_fixed_projection_block_decomp_isometry}
    \lean{MPSTensor.exists_blockDecomp_of_adjoint_fixed_projections_with_letter_and_isometry}
    \leanok
    \uses{thm:commuting_projection_block_decomp_isometry}
    Let $A$ be left-canonical and let $(P_k)$ be orthogonal projections
    summing to $\Id$ and fixed by the adjoint transfer map. Then the same
    sector decomposition and support isometries of
    Theorem~\ref{thm:commuting_projection_block_decomp_isometry} exist. In
    particular, the adjoint fixed-point condition $\E_A^\dagger(P_k)=P_k$
    gives
    \[
        P_k A^i=A^i P_k,
    \]
    which gives the corner-letter identity
    \[
        \varphi_k(C_k^i)=P_k A^i P_k.
    \]
    \begin{proof}\leanok
        \uses{thm:commuting_projection_block_decomp_isometry}
        The fixed-point equation $\E_A^\dagger(P_k)=P_k$, together with
        left-canonicality $\sum_i (A^i)^\dagger A^i=\Id$, places $P_k$ in the
        multiplicative domain of $\E_A^\dagger$, forcing
        \[
            P_k A^i=A^i P_k
        \]
        for every letter $A^i$.  The result then follows from
        Theorem~\ref{thm:commuting_projection_block_decomp_isometry}.
    \end{proof}
\end{theorem}

\begin{theorem}[Spectral cyclic sectors with support isometries]%
    \label{thm:cyclic_sector_decomp_after_blocking_isometry}
    \lean{MPSTensor.exists_cyclic_sector_decomp_after_blocking_with_letter_and_isometry}
    \leanok
    \uses{thm:peripheral_cyclic_structure,
        thm:adjoint_fixed_projection_block_decomp_isometry}
    Let $A$ be a left-canonical irreducible tensor, and suppose the adjoint
    transfer map $\E_{A^\dagger}$ has a positive-definite fixed point, is
    irreducible as a positive map, and has peripheral spectrum
    \[
        \sigma_\partial(\E_{A^\dagger})
        =
        \{\gamma^j : j\in \Z/m\Z\}
    \]
    for a primitive $m$th root of unity $\gamma$, with $m\ge 1$. Then the
    spectral cyclic-sector decomposition of $A^{[m]}$ has compression maps
    which may be chosen with rectangular support isometries $V_k$ satisfying
    \[
        V_k^\dagger V_k=\Id,\qquad V_kV_k^\dagger=P_k,\qquad
        \varphi_k(X)=V_kXV_k^\dagger .
    \]
    These isometries also realize the blocked corner letters
    $\varphi_k(C_k^\alpha)=P_k(A^{[m]})^\alpha P_k$.
    \begin{proof}\leanok
        \uses{thm:peripheral_cyclic_structure,
            thm:adjoint_fixed_projection_block_decomp_isometry}
        The cyclic peripheral-spectrum decomposition applied to
        $\E_{A^\dagger}$ provides projections $P_k$ satisfying
        $\E_{A^\dagger}(P_{k+1})=P_k$.  Iterating $m$ times gives
        $(\E_{A^\dagger})^m(P_k)=P_k$, and the identity
        $(\E_{A^\dagger})^m=\E_{(A^{[m]})^\dagger}$ shows that each $P_k$ is
        fixed by the adjoint transfer map of $A^{[m]}$. Applying
        Theorem~\ref{thm:adjoint_fixed_projection_block_decomp_isometry} to
        that blocked tensor gives the compressed sector tensors and the
        support isometries.
    \end{proof}
\end{theorem}

\subsection{Reduction to primitive blocks}%
\label{sec:reduction_to_normal_form}

This subsection completes the reduction after period removal. A
TP-primitive-irreducible block is normal, and an arbitrary tensor becomes, after a
positive blocking length, a weighted family of left-canonical primitive blocks
whose total bond dimension is bounded by the original bond dimension. For two
tensors with the same MPV family, the cyclic sectors can be expressed over one
common blocked alphabet, which is the form used by the Fundamental Theorem.

\begin{theorem}[TP-primitive irreducible tensor is normal]%
    \label{thm:normal_tp_primitive_irred}
    \lean{MPSTensor.isNormal_of_tp_primitive_irreducible}
    \leanok
    \uses{def:primitive_mps, def:irreducible_tensor, def:normal}
    Let $A$ be a left-canonical MPS tensor with $D > 0$, irreducible transfer
    map, and peripheral spectrum $\{1\}$. Then $A$ is normal.
\end{theorem}

\begin{proof}\leanok
    \uses{thm:normal_from_primitive_posdef, thm:psd_fp_pd_irred}
    Irreducibility upgrades the PSD fixed point to positive definite
    (Theorem~\ref{thm:psd_fp_pd_irred}), and
    Theorem~\ref{thm:normal_from_primitive_posdef} then gives normality.
\end{proof}

An arbitrary translation-invariant tensor reduces, after blocking, to a weighted
direct sum of primitive irreducible blocks.

\begin{theorem}[Arbitrary tensor to primitive block decomposition]%
    \label{thm:tp_primitive_blockdecomp}
    \lean{MPSTensor.exists_tp_primitive_blockDecomp_after_blocking}
    \leanok
    \uses{def:blocked_tensor, def:same_mpv2_pos}
    For any MPS tensor $A$, there exist $p\ge 1$, nonzero weights
    $(\mu_k)_{k=1}^r$, and blocks $(B_k)_{k=1}^r$ such
    that, for every blocked $\sigma$ of positive length $N$,
    \[
        V^{(N)}\!(A^{[p]})_\sigma
        = V^{(N)}\!(\textstyle\bigoplus_{k=1}^{r} \mu_k B_k)_\sigma.
    \]
    The retained block dimensions satisfy $\sum_{k=1}^{r}D_k\le D$. Each block
    is left-canonical and primitive,
    \[
        \sum_i (B_k^i)^\dagger B_k^i = \Id,
        \qquad
        \sigma_\partial(\E_{B_k})=\{1\},
    \]
    and each bond dimension is positive.
\end{theorem}

\begin{proof}\leanok
    \uses{thm:tp_gauge_arbitrary, thm:common_blocking_period,
        lem:block_weights_ne_zero,
        thm:same_mpv2_pos_block_tensor_to_tensor_from_blocks}
    Theorem~\ref{thm:tp_gauge_arbitrary} supplies the left-canonical nontrivial
    blocks, the positive-length equality, and the bound $\sum_kD_k\le D$.
    Theorem~\ref{thm:common_blocking_period} provides a common blocking period
    $p$ making all transfer maps primitive,
    Lemma~\ref{lem:block_weights_ne_zero} keeps the blocked weights nonzero, and
    Lemma~\ref{thm:same_mpv2_pos_block_tensor_to_tensor_from_blocks} carries
    the positive-length equality through blocking.
\end{proof}

\begin{theorem}[Unconditional common primitive block decompositions]%
    \label{thm:unconditional_common_primitive_irreducible_blocks}
    \lean{MPSTensor.unconditional_commonPrimitiveIrreducibleBlocks}
    \leanok
    \uses{thm:ft_after_blocking_common_length_common_sector_theorem,
        thm:common_blocked_cyclic_sector_reindexed_nonzero_part,
        def:same_mpv2_pos}
    Let $A$ and $B$ have the same MPV family at every positive length. Then there is a positive blocking
    length at which both blocked tensors have the same positive-length MPV
    family as weighted direct sums of trace-preserving, primitive,
    tensor-irreducible blocks with positive bond dimensions and nonzero weights.
    The two weighted nonzero-sector families have the same positive-length MPV
    family.
\end{theorem}

\begin{proof}
    \leanok
    \uses{thm:ft_after_blocking_common_length_common_sector_theorem,
        thm:common_blocked_cyclic_sector_reindexed_nonzero_part}
    Apply
    Theorem~\ref{thm:ft_after_blocking_common_length_common_sector_theorem} to
    choose the common blocking length and the two common cyclic-sector families.
    For either one of these families, write $G_k$ for its cyclic-sector blocks,
    and write $(\nu_x,C_x)$ for the weights and blocks of the flattened
    common-sector family from
    Lemma~\ref{thm:common_blocked_cyclic_sector_reindexed_nonzero_part}. The
    common-alphabet nonzero-part identity gives, for every length $N$ and
    blocked word $\sigma$,
    \[
        V^{(N)}\!(\textstyle\bigoplus_k \mu_k^p G_k^{[p]})_\sigma
        =
        V^{(N)}\!(\textstyle\bigoplus_x \nu_x C_x)_\sigma ,
    \]
    Applying this on the two sides of the comparison yields a positive length
    $p$ and weighted primitive block families such that, for every $N\ge 1$ and
    blocked word $\sigma$,
    \[
        V^{(N)}\!(A^{[p]})_\sigma
        =
        V^{(N)}\!(\textstyle\bigoplus_x \mu^A_x A_x)_\sigma,
        \qquad
        V^{(N)}\!(B^{[p]})_\sigma
        =
        V^{(N)}\!(\textstyle\bigoplus_y \mu^B_y B_y)_\sigma,
    \]
    with the two weighted nonzero parts equal at every positive length. The
    structural properties and the nonvanishing of the weights come from the
    common-sector families.
\end{proof}

A weighted block sum with distinct sectors is the simplest sector decomposition:
one block per sector, each of multiplicity one. Equal-modulus and repeated-copy
sectors are not single scalar powers, already for $A = C \oplus (-C)$ with
$V^{(N)}(A)_\sigma = (1 + (-1)^N)\,V^{(N)}(C)_\sigma$; they are compared by the
basis-of-normal-tensors construction of Chapter~\ref{ch:bnt}.

\begin{definition}[Single-copy sector decomposition]
    \label{def:trivial_sector_decomp}
    \lean{MPSTensor.trivialSectorDecomp}
    \leanok
    \uses{def:sector_weight_data}
    Given nonzero weights $(\mu_k)_{k=1}^r$ and blocks $(B_k)_{k=1}^r$, the
    \emph{single-copy sector decomposition} has $r$ sectors, one basis block
    $B_k$ per sector, multiplicity~$1$, and sector weight $\mu_k$. As a finite
    sum,
    \[
        \sum_{k=1}^{r} \mu_k^{\,N} V^{(N)}(B_k)_\sigma
        =  V^{(N)}\!(\textstyle\bigoplus_k \mu_k B_k)_\sigma.
    \]
\end{definition}

\begin{lemma}[Single-copy sector decomposition preserves the represented family]
    \label{lem:trivial_sector_decomp_same_mpv}
    \lean{MPSTensor.sameMPV₂_trivialSectorDecomp}
    \leanok
    \uses{def:trivial_sector_decomp}
    The single-copy sector decomposition of $(\mu_k,B_k)_k$ represents the same
    MPS family as $\bigoplus_k \mu_k B_k$.
\end{lemma}

\begin{proof}\leanok
    \uses{def:trivial_sector_decomp}
    Expanding the sector decomposition gives $\sum_k \mu_k^N V^{(N)}(B_k)_\sigma$,
    the finite block-sum expression for $\bigoplus_k \mu_k B_k$.
\end{proof}


\section{Further weighted-block identities}

\begin{lemma}[Canonical form from peripheral primitivity]%
    \label{thm:canonical_from_primitive}
    \lean{MPSTensor.IsCanonicalForm.of_peripheral_primitive}
    \leanok
    \uses{def:is_canonical_form, def:injective, def:tp_kraus,
        def:primitive_channel}
    Let $(\mu_k,A_k)_{k=1}^r$ be nonzero weights and injective left-canonical
    blocks with positive bond dimensions, non-increasing moduli
    $|\mu_1|\ge\cdots\ge|\mu_r|>0$, and $\sigma_\partial(\E_{A_k}) = \{1\}$. Then
    $(\mu_k,A_k)_{k=1}^r$ satisfies the canonical-form conditions of
    Definition~\ref{def:is_canonical_form}, and $O_{A_k A_k}(N) \to 1$ for every
    $k$.
\end{lemma}

\begin{proof}
    \leanok
    \uses{def:is_canonical_form, thm:primitive_compl_lt_one, thm:primitive_overlap}
    The injectivity, left-canonical equation, weight ordering, and nonzero-weight
    clauses are exactly the hypotheses.  For each block, peripheral primitivity
    gives the complementary spectral gap
    $\rho_{\spec}(\E_{A_k}-P) < 1$ by Theorem~\ref{thm:primitive_compl_lt_one}.
    Theorem~\ref{thm:primitive_overlap} then yields $O_{A_k A_k}(N)\to 1$.
\end{proof}

\begin{lemma}[Common-alphabet flattening]%
    \label{thm:common_blocked_cyclic_sector_reindexed_samempv}
    \lean{MPSTensor.CommonBlockedCyclicSectorFamily.reindexed_sameMPV₂}\leanok
    \uses{thm:exists_common_blocked_cyclic_sector_family}
    With $p = m_k e_k$ for each $k$, the iterated blocked tensor
    $(B_k^{[m_k]})^{[e_k]}$ agrees with $B_k^{[p]}$ in MPV family, after the
    canonical identification of the $p$-blocked alphabet $d^{[p]}$ with the
    $e_k$-fold tensor power of the $m_k$-blocked alphabet $d^{[m_k]}$.
\end{lemma}

\begin{lemma}[Reindexed nonzero-part identity]%
    \label{thm:common_blocked_cyclic_sector_reindexed_nonzero_part}
    \lean{MPSTensor.CommonBlockedCyclicSectorFamily.reindexed_nonzero_part}\leanok
    \uses{thm:common_blocked_cyclic_sector_reindexed_samempv,
        thm:common_blocked_cyclic_sector_block_tensor_common_reindexed}
    The powered nonzero-part decomposition
    $\bigoplus_k \mu_k^p B_k^{[p]}$ carries over under the common-alphabet
    identification of
    Lemma~\ref{thm:common_blocked_cyclic_sector_reindexed_samempv}: the
    common-alphabet sectors have the same nonzero-part decomposition as the direct
    $p$-blocked tensors.
\end{lemma}

\begin{lemma}[Direct blocking in the common alphabet]%
    \label{thm:common_blocked_cyclic_sector_block_tensor_common_reindexed}
    \lean{MPSTensor.CommonBlockedCyclicSectorFamily.blockTensor_sameMPV₂_commonReindexedBlock}\leanok
    For every original block $B_k$, direct blocking at the common length
    $p=m_k e_k$ agrees, as an MPV family, with the same block transported to the
    common blocked alphabet. Write $\widetilde B_k$ for this common-alphabet
    image. Then
    \[
        V^{(N)}\!(B_k^{[p]})_\sigma
        =
        V^{(N)}\!(\widetilde B_k)_\sigma.
    \]
\end{lemma}
